\subsection{极化}\label{gq}

实践中，量子化通常通过把相空间变量分成坐标与动量来进行。然后定义由坐标的函数构成的 Hilbert 空间。我们在 $\R^{2n}$ 的情形已经讨论过这一点。在几何量子化 \cite{Ko,So,Wo,Snia,Hall} 中，这一过程被形式化为“极化”的选取。虽然还有更一般的可能，但两个基本情形，也是本文主要考虑的情形，是\textbf{实极化}与\textbf{复极化}。

在实极化最简单的形式中，把 $M$ 等同于某个 $N$ 的余切丛 $T^*N$ 的全空间。假定 $M$ 的辛结构就是 $T^*N$ 作为余切丛的自然辛结构，或者是它与 $N$ 上某个闭 2-形式之拉回的和。用 $N$ 上的局部坐标 $q^i$，$i=1,\cdots,n$ 来说，这意味着辛形式为 $\omega =\sum_i \d p_i \d q^i +\frac{1}{2}\sum_{i,j}\alpha_{ij}(q) \d q^i\d q^j$，其中 $p_i$ 是 $\pi:T^*N\to N$ 的纤维上合适的线性函数，$\alpha=\frac{1}{2}\sum_{i,j}\alpha_{ij}(q) \d q^i\d q^j$ 是 $N$ 上的闭 2-形式。更简要地说，$\omega=\sum_i \d p_i \d q^i+\pi^*(\alpha)$，其中 $\pi:M=T^*N\to N$ 是自然投影。在此情形下，$q^i$ 称为坐标，$p_i$ 称为动量。余切丛的纤维称为极化的\textbf{叶}。

定义 Hilbert 空间 $\H$，它由 $N$ 上取值于预量子线丛 $\frL$ 的半密度\footnote{在实流形 $N$ 上的局部坐标系 $\vec q=(q^1,q^2,\cdots, q^n)$ 中，一个\textbf{密度}就是实值函数 $\sigma(\vec q)$。在局部坐标变换下，$\sigma$ 的变换方式使得测度 $|\d q^1\d q^2 \cdots \d q^n |\sigma(\vec q)$ 不变。\textbf{半密度}是对象 $\psi(\vec q)$，使得 $\sigma=\psi^2$ 是密度。因此，在给定坐标系中 $\psi(\vec q)$ 就是函数，而在局部坐标变换下，$\psi(\vec q)$ 的变换方式使得测度 $|\d q^1 \d q^2 \cdots \d q^n| \,\psi(\vec q)^2$ 不变。存在平凡实线丛 $\K\to N$，其截面是（实）密度；$\K$ 有平方根 $\K^{1/2}$，也是平凡的，其截面是半密度。我们一般用 $\K$ 表示实流形 $N$ 上的实密度丛，用 $K$ 表示复流形 $X$ 的典则丛，用 $\bar K$ 表示 $K$ 的复共轭（等价地，是把 $X$ 看成具有相反复结构的复流形时的典则丛）。若把复流形 $X$ 视为实流形，则其密度丛 $\K$ 与典则丛 $K$ 由 $\K=K\otimes \bar K=|K|^2$ 相联系。}构成。换言之，若 $\K^{1/2}$ 是 $N$ 上以半密度为截面的实线丛，则 $\H$ 中的向量由 $\K^{1/2}\otimes \frL\to N$ 的截面 $\psi$ 表示，内积为 $\la\psi,\chi\ra=\int_N \bar\psi \chi$。这个定义有时重述如下：先把 $\K^{1/2}\to N$ 拉回到 $T^*N$ 上的实线丛，仍记作 $\K^{1/2}$；于是 $\H$ 中的向量可描述为 $\K^{1/2}\otimes \frL$ 在 $T^*N$ 上的截面，它在限制到 $T^*N\to N$ 的任一纤维上时协变常数。这是有意义的，因为（其曲率为 $\omega$）$\frL$ 在限制到极化的任一叶上时是平坦的，而 $\K^{1/2}$ 作为从 $N$ 的拉回，沿每个叶也有自然的平坦联络。Hilbert 空间内积于是可写成积分 $\la\psi,\chi\ra=\int_\sigma \bar\psi\chi$，其中 $\sigma$ 是 $T^*N\to N$ 的任一截面。

现在讨论复极化。$M$ 的\textbf{复极化}定义为 $M$ 上复结构 $J$ 的选取，使得辛形式 $\omega$ 为 $(1,1)$ 型，从而预量子线丛 $\frL$ 成为全纯线丛。若 $\omega$ 关于 $J$ 既为 $(1,1)$ 型又为正，因而代表 $M$ 上的 Kähler 形式，则称该极化为 \textbf{Kähler 极化}。这是实践中最常考虑的情形，尽管我们并不想局限于此。

设 $K$ 是 $M$ 的典则线丛，$K^{1/2}$ 是 $K$ 的平方根，暂假定这样的平方根存在。Kähler 量子化中通常考虑 $\frL$ “充分丰沛”的情形，使得 $H^r(M,K^{1/2}\otimes \frL)=0$（$r>0$）。（在充分接近经典极限时，即 $\omega$ 的周期充分大时，这总是成立的。）于是定义全纯截面构成的 Hilbert 空间\footnote{在实极化的情形，实半密度丛 $\K^{1/2}$ 是量子态定义的一部分，否则 Hilbert 空间内积的定义 $\la\psi,\chi\ra=\int_N\bar\psi\chi$ 就没有意义。在 Kähler 极化的情形，典则丛 $K$ 有来自 $M$ 的 Kähler 度量的 Hermite 度量，所以无论定义中是否包含 $K^{1/2}$ 因子，都可以定义 Hilbert 空间内积（\ref{cph}）。通常包含这个因子有两个理由：(1) 这样可以对不同类型的极化给出更统一的几何量子化描述；(2) 在某些情形下它导致更自然的结果，例如用相空间的 Kähler 极化量子化谐振子时基态能量的计算。参见例如 \cite{Hall} 第 23.7 节。}
\be\label{cph}\H=H^0(M,K^{1/2}\otimes \frL),\ee
其 Hilbert 空间内积为
\be\label{inpro}\la s,s'\ra=\int_M \d\mu \, \bar s s',\ee
其中 $\d\mu$ 是 $M$ 上的辛测度，而 $K$ 上的 Kähler 度量与 $\frL$ 的酉结构把 $\bar s s'$ 变成实值函数。若 $\omega$ 不是 Kähler 型或 $\frL$ 不是充分丰沛，其推广见第 \ref{branecp} 节。

上一段中已经隐含了：几何量子化存在反常，它与所谓的\textbf{亚辛校正}有关。一般地，用于定义 Hilbert 空间的对象 $K^{1/2}$ 并不作为复线丛存在；即使存在，也未必唯一。若 $K^{1/2}$ 存在，它定义了 $M$ 上的一个 spin 结构；反之，spin 结构决定一个 $K^{1/2}$ 的选取。$M$ 容许 spin 结构的障碍是 $w_2(M)$，即 $M$ 的第二 Stiefel--Whitney 类。然而，在任何复流形上，$K^{1/2}$ 总作为 $\spinc$ 结构典则地存在。一般的 $\spinc$ 结构是定义作用在“荷为 1”的 spin 1/2 场上的 Dirac 算符所需的数据。在 $n$ 维复流形上，对应于典则 $\spinc$ 结构 $K^{1/2}$ 的 Dirac 算符，就是作用在 $\oplus_{r=0}^n\Omega^{0,r}(M)$ 上的 $\bar\partial$ 算符，其中 $\Omega^{0,r}(M)$ 是 $(0,r)$-形式丛。由于 spin 的障碍是 2 阶的，两个 $\spinc$ 结构的张量积是普通复线丛。因此，若几何量子化的“预量子线丛 $\frL$”实际上是 $\spinc$ 结构，则 $K^{1/2}\otimes \frL$ 是普通复线丛，量子 Hilbert 空间的定义 \eqref{cph} 就有意义。

更一般地，若 $\SS$ 是流形 $M$ 上任一 $\spinc$ 结构（这里不假定 $M$ 是复流形，也不假定 $\SS$ 与预量子线丛有关），则 $\SS^2$ 是普通线丛。此线丛有第一 Chern 类 $x=c_1(\SS^2)$，它是 $w_2(M)$ 的整数提升，即 $x\in H^2(M,\Z)$ 模 2 约化后是 $w_2(M)\in H^2(M,\Z_2)$。若 $\SS$ 是带联络的 $\spinc$ 结构，则普通线丛 $\SS^2$ 有联络与曲率；我们把 $\SS^2$ 曲率的一半定义为 $\SS$ 的曲率。若此曲率为零，则称 $\spinc$ 结构\textbf{平坦}。按此定义，$\SS$ 容许平坦联络的条件就是 $x$ 应当是挠元；换言之，若 $w_2(M)$ 有一个挠的整数提升，则 $M$ 有平坦 $\spinc$ 结构。

按这些定义，Kähler 流形 $M$（Kähler 形式为 $\omega$）上的“预量子线丛” $\frL$ 实际上应定义为 $M$ 上曲率为 $\omega$ 的 $\spinc$ 结构。

一旦把几何量子化的“预量子线丛”重新解释为 $\spinc$ 结构，就应问：这会不会给余切丛 $M=T^*N$ 的量子化讨论带来困难？事实上，这一反常不影响 $T^*N$ 上预量子线丛的分类，因为 $T^*N$ 总有典则的平坦 $\spinc$ 结构 $\SS_0$，因此 $T^*N$ 上曲率为 $\omega$ 的 $\spinc$ 结构就是 $\SS_0\otimes \L$，其中 $\L$ 是曲率为 $\omega$ 的酉联络的普通复线丛。

为解释这一点，先要从拓扑上理解为什么 $T^*N$ 容许平坦 $\spinc$ 结构。$T^*N$ 的切丛在拓扑上是 $TN\oplus TN$（$N$ 的切丛的两个拷贝）到 $T^*N$ 的拉回。Whitney 和公式给出 $w_2(TN\oplus TN)= w_1(TN)^2$，其中 $w_1(TN)$（通常简记作 $w_1(N)$）是 $N$ 可定向性的障碍。所以 $w_2(T^*N)$ 是 $w_1(N)^2$ 到 $T^*N$ 的拉回。

设 $\varepsilon=\det \,TN$ 是 $N$ 的定向丛，一个实线丛。则 $w_1(N)=w_1(\varepsilon)$。对象 $\h\varepsilon=\varepsilon\oplus \varepsilon\cong \varepsilon\otimes_\R\C$ 是 $N$ 上的复线丛。其第一 Chern 类 $x=c_1(\h\varepsilon)\in H^2(N,\Z)$ 模 2 约化为 $w_2(\h\varepsilon)=w_1(\varepsilon)^2=w_2(T^*N)\in H^2(N,\Z_2)$。$x$ 当然是挠类，因为 $\h\varepsilon$ 作为实线丛的直和，容许结构群有限的平坦联络。所以 $x$ 是 $w_2(T^*N)$ 的整数提升，而且是挠的，因此 $T^*N$ 容许平坦 $\spinc$ 结构。

为显式构造 $T^*N$ 上的平坦 $\spinc$ 结构，我们先在丛 $TN\oplus TN\to N$ 上构造平坦 $\spinc$ 结构；它到 $T^*N$ 的拉回就是所需的 $T^*N$ 上的平坦 $\spinc$ 结构。取 $TN$ 的局部正交标架，引入两组 gamma 矩阵 $\psi_k, \chi_k$，$k=1,\cdots,n$，满足 $\{\psi_k,\psi_l\}=\{\chi_k,\chi_l\}=2\delta_{kl}$，$\{\psi_k,\chi_l\}=0$。$TN\oplus TN$ 上的 $\spinc$ 结构是此 Clifford 代数在其上不可约作用的向量丛。引入湮灭与产生算符 $a_k=(\psi_k+\i \chi_k)/2,$ $a^\dagger_k=(\psi_k-\i\chi_k)/2$，于是 $\{a_k,a^\dagger_l\}=\delta_{kl}$，$\{a_k,a_l\}=\{a^\dagger_k,a^\dagger_l\}=0$。局部地，可以从被 $a_k$ 零化的态 $|\negthinspace\downarrow\ra$ 出发，添加态 $a^\dagger_{k_1}a^\dagger_{k_2}\cdots a^\dagger_{k_r}|\negthinspace\downarrow\ra$（$r=1,\cdots, n$），构造 Clifford 代数的模。特别地，态 $|\negthinspace\uparrow\ra=a_1^\dagger a_2^\dagger\cdots a_n^\dagger|\negthinspace\downarrow\ra$ 被 $a^\dagger_k$ 零化。若 $N$ 可定向，此构造给出 $T^*N$ 上典则的 spin 结构。若 $N$ 不可定向，则两个态 $|\negthinspace\downarrow\ra$ 与 $|\negthinspace\uparrow\ra$ 整体上不对称；若其中之一是 $N$ 上平凡丛的截面，则另一个是 $\h\varepsilon$ 的截面。这种不对称性意味着\footnote{$TN\oplus TN$ 的结构群是 $\SO(2n)$，相应的 spin 丛的结构群是 $\Spin(2n)$。这样的 spin 丛有非退化的 $\Spin(2n)$-不变二次型（依 $n$ 为对称或反对称）。此不变形式限制为 $|\negthinspace\uparrow\ra \otimes |\negthinspace\downarrow\ra\to \C$ 上的非退化配对。若 $|\negthinspace\downarrow\ra$ 是线丛 $\ell$ 的截面，则 $|\negthinspace\uparrow\ra$ 是 $\ell\otimes \varepsilon$ 的截面，于是配对的存在意味着 $\ell^2\otimes \varepsilon$ 平凡，从而 $2 c_1(\ell)+x=0$。这意味着 $x$ 的模 2 约化为零；既然 $x$ 的模 2 约化是 $w_2(T^*N)$，我们就如所预期地得知：文中的构造只有在 $w_2(T^*N)=0$ 时才给出 spin 丛，否则给出平坦 $\spinc$ 丛。}我们构造出的其实是 $\spinc$ 丛而非 spin 丛；不过，由于 $\h\varepsilon$ 有自然的平坦结构，这个 $\spinc$ 结构是平坦的。所以一般结论是：$T^*N$ 总有典则的平坦 $\spinc$ 结构。

几何量子化的主要局限在于：若 $M$ 存在一个极化，则存在许多不同的极化，因为对称群 $\G$ 的任何一般元素都会把极化 $\P$ 映到不同的极化 $\P'$。在合理条件下，可以用不同极化定义的 Hilbert 空间之间相对自然的映射。\footnote{\label{BKS}这是用几何量子化的 Blattner--Kostant--Sternberg 核做到的。对两个适当横截的极化，定义最为直接。简要介绍见 \cite{Hall} 第 23.8 节，更多细节见 \cite{Snia}。另见下面的脚注 \ref{BKStwo}。}但这些映射并不满足人们为建立不同量子化之间的等价所希望的条件。人们希望的是如下条件：给定导致 Hilbert 空间 $\H_i$ 的极化 $\P_i$，希望有一族酉映射 $\Phi_{ij}:\H_i\to \H_j$，满足 $\Phi_{ii}=1$、$\Phi_{ij}\Phi_{ji}=1$（从而 $\Phi_{ij}$ 是同构）以及 $\Phi_{ij}\Phi_{jk}=\Phi_{ik}$（从而不同 Hilbert 空间之间的这些同构是相容的）。但这些条件一般并不满足；这一陈述是经典到量子物理过渡中反常的另一种表达。

不同 Hilbert 空间之间相容性的条件可以稍加放宽，允许 $\Phi_{ij}\Phi_{jk}=c_{ijk} \Phi_{ik}$，其中 $c_{ijk}$ 是取值于 $\UU(1)$ 的中心因子。此时人们会说这些 $\H_i$ “投影等价”，意即它们之间存在自然定义到模 1 标量因子的等价。这个稍宽的极化间等价条件一般也不满足，因为经典力学到量子力学的反常并不只涉及一个中心扩张。

关于不同极化何时导致等价量子化，没有一般理论，但确实有许多重要情形会发生。最熟悉、最重要的例子出现在量子化 $M=\R^{2n}$ 时。此时，如已指出，一旦给定仿射线性结构，用不同线性极化定义的 Hilbert 空间之间确实有投影等价。$\R^{2n}$ 的\textbf{线性极化}，在实情形下是把 $\R^{2n}$ 上的线性函数分成坐标与动量，在复情形下是选择 $\R^{2n}$ 上哪些复值线性函数应被视为“全纯的”。本文中我们还会看到几个额外例子，其中不同极化导致至少投影等价的量子化。每个这样的情形都会导致重要的结果。

几何量子化还有一个更特殊的局限：存在不容许任何极化的辛流形 \cite{Gotay}。这样的辛流形能否以某种合理的方式量子化，至少在某些情形下？这个问题是寻找几何量子化替代方案的一个动机。

\begin{yizhu}[“亚辛校正”与 $\Spinc$ 结构]
几何量子化里这个 $K^{1/2}$ 因子，文献中常称 \textbf{metaplectic correction}（亚辛校正）。它并非可有可无的装饰：若要求量子化在不同极化之间协变，半形式的权重必须取对。原文的要点是，$K^{1/2}$ 作为线丛未必存在（障碍是 $w_2(M)$ 的非零性），但作为 $\Spinc$ 结构却总存在——复流形上 $\bar\partial$ 作用在 $(0,\bullet)$-形式丛上就给出这个典则 $\Spinc$ 结构，其“平方”恰是 $K$。因此原文把预量子线丛改述为“曲率为 $\omega$ 的 $\Spinc$ 结构”；在 Kähler 情形，这与“线丛张量 $K^{1/2}$”的描述一致。[来源：I 级]
\end{yizhu}
